\chapter{Block Diagram}
\label{ch:block}

\section{Purpose}

The block diagram shows the components of the system, the boundary that
separates them, and what flows across it. It records the two decisions that
constrain everything else in this document: the split between code running
inside GENESYS and code running in a separate process, and the confinement of
vendor assemblies to a single component.

\umlfigure{block-diagram}{System block diagram}{fig:block}

\section{Elements}

The diagram is divided into two groups. \emph{In GENESYS} holds the components
that run inside the host application. \emph{.NET Process} holds those that run
separately. The GENESYS Model and the Code Repository sit outside both, because
the tool reads and writes them but does not implement either.

\begin{longtable}{L{0.22\textwidth}L{0.70\textwidth}}
	\caption{Block diagram elements}\label{tab:block-elements}\\
	\toprule
	\textbf{Component} & \textbf{Responsibility} \\
	\midrule
	\endfirsthead
	\toprule
	\textbf{Component} & \textbf{Responsibility} \\
	\midrule
	\endhead

	Repository Client & Connects to GENESYS, reads the project schema, and creates
	or updates entities. The only component that references the vendor client
	assemblies. \\
	\addlinespace

	Model Mapper & Translates between the intermediate representation and GENESYS
	entities, including states, transitions and behavioural constructs. Sends and
	receives that representation across the process boundary. \\
	\addlinespace

	Code Extractor & Loads a solution, resolves symbols through the compiler's
	semantic model, and produces control flow graphs. Reads from the Code
	Repository. \\
	\addlinespace

	Abstraction & Applies significance rules in both directions, discarding
	constructs that do not belong at the level a systems engineer reads. Sits
	between the extractor and the mapper, and between the mapper and the
	generator. \\
	\addlinespace

	Code Generator & Emits source implementing a model's control flow, and writes
	it to the Code Repository. \\
	\addlinespace

	GENESYS Model & External. The repository holding the system model. \\
	\addlinespace

	Code Repository & External. Holds the source files the tool reads and writes. \\
	\bottomrule
\end{longtable}

\section{Data Flow}

Two paths run through the diagram in opposite directions. Code to model begins
at the Code Repository, passes through the Code Extractor and the Abstraction
component, crosses the process boundary into the Model Mapper, and reaches the
GENESYS Model through the Repository Client. Model to code reverses that route,
leaving through the Code Generator.

The Abstraction component is the junction of both paths, which is why it sits
between the extractor and the generator rather than inside either.

\section{Design Decisions}

\paragraph{Process boundary.}
The vendor client assemblies target .NET Framework 4.x, while the compiler
platform used for extraction targets a current .NET runtime. Loading both into
one process risks assembly version conflicts, and a fault in the parser would
otherwise terminate the host application. This separation cannot be retrofitted
later without redesign.

\paragraph{Confinement of vendor assemblies.}
Only the Repository Client references the vendor libraries. Every other
component builds and runs with no licence and no GENESYS installation, which
matters for a team with a limited number of licence activations.

\paragraph{A single abstraction component serving both directions.}
Rather than filtering separately on each path, one component holds the
significance rules. The rules are therefore stated once and applied
consistently, whichever direction the translation runs.

\paragraph{External systems marked explicitly.}
The GENESYS Model and the Code Repository are drawn outside the process groups
to make clear that neither is built by this project, and that the tool performs
no version control operations of its own.
